\chapter{System Features}

\section{Feature 1: Prism AI Toolbar Button}

\subsection{Description and Priority}
The Prism AI Button is a custom toolbar component integrated into the Chromium browser toolbar. It serves as the primary entry point for all AI-assisted features. It is a \textbf{high-priority} feature as it ties the browser UI to the AI backend.

\subsection{Implementation Details}
The button is implemented in C++ as the class \texttt{PrismAIButton}, which extends Chromium's \texttt{ToolbarButton} class. Key implementation aspects:
\begin{itemize}[leftmargin=*]
    \item The class is defined in \texttt{chrome/browser/ui/views/toolbar/prism\_ai\_button.h} and implemented in \texttt{prism\_ai\_button.cc}.
    \item It uses the custom vector icon \texttt{kPrismAiIcon} defined in \texttt{chrome/app/vector\_icons/prism\_ai.icon}, which renders a prism-shaped double triangle on a 16$\times$16 canvas.
    \item The tooltip text is set to ``Prism AI''.
    \item When clicked, \texttt{OnButtonPressed()} opens a new foreground tab to \texttt{http://127.0.0.1:7788} using the browser's \texttt{OpenURL()} method with \texttt{PAGE\_TRANSITION\_TYPED} transition type.
    \item The button is added as a child view in \texttt{ToolbarView::Init()} after the Home button, and a raw pointer \texttt{prism\_ai\_button\_} is maintained in \texttt{ToolbarView}.
\end{itemize}

\subsection{Functional Requirements Supported}
FR-01

\section{Feature 2: Always-On Chrome DevTools Protocol (CDP)}

\subsection{Description and Priority}
CDP is the communication backbone between the AI agent and the browser. Without CDP being active, the AI agent cannot control or inspect browser state. This feature is \textbf{high-priority}.

\subsection{Implementation Details}
In \texttt{chrome/app/chrome\_main.cc}, the following logic is injected at browser startup:
\begin{itemize}[leftmargin=*]
    \item After the command line is initialised, the code checks whether \texttt{--remote-debugging-port} or \texttt{--remote-debugging-pipe} switches are already set.
    \item If neither is present, \texttt{--remote-debugging-port=9222} is appended programmatically, guaranteeing CDP is always available.
    \item This eliminates the need for users to start the browser with special flags.
\end{itemize}
The CDP endpoint is accessible at \texttt{ws://127.0.0.1:9222} from the local machine only, providing a WebSocket interface for tab management, DOM access, JavaScript execution, and input simulation.

\subsection{Functional Requirements Supported}
FR-02

\section{Feature 3: AI Agent Backend}

\subsection{Description and Priority}
The AI agent server is the intelligence core of Prism. It is a Python application that serves the chat UI, processes user commands using an LLM, and drives the browser via CDP. This is a \textbf{high-priority} feature.

\subsection{Implementation Details}
\begin{itemize}[leftmargin=*]
    \item The server listens on \texttt{http://127.0.0.1:7788} (loopback only).
    \item On receiving a user command, the server sends the command to a locally-running LLM (e.g., via the Ollama REST API at \texttt{http://127.0.0.1:11434}).
    \item The LLM decomposes the command into a sequence of browser automation actions (a ``plan'').
    \item Each action is executed by sending the corresponding CDP command over WebSocket to the browser at \texttt{ws://127.0.0.1:9222}.
    \item Common CDP actions include: \texttt{Target.createTarget}, \texttt{Page.navigate}, \texttt{DOM.getDocument}, \texttt{Input.dispatchMouseEvent}, and \texttt{Runtime.evaluate}.
    \item Results and status messages are returned to the web UI in real time.
\end{itemize}

\subsection{Functional Requirements Supported}
FR-03, FR-05, FR-06, FR-07, FR-08

\section{Feature 4: Voice-to-Text Input}

\subsection{Description and Priority}
Voice command support makes Prism AI Browser accessible to users with limited mobility and enables hands-free productivity. This is a \textbf{medium-priority} feature.

\subsection{Implementation Details}
\begin{itemize}[leftmargin=*]
    \item The AI agent web UI provides a microphone button that activates the browser's Web Speech API or a locally hosted STT engine.
    \item The transcribed speech is populated in the text command field.
    \item The user can review the transcribed text before submitting or it can be submitted automatically.
    \item All audio processing is performed locally to maintain privacy compliance.
\end{itemize}

\subsection{Functional Requirements Supported}
FR-04

\section{Feature 5: Content Summarisation and Q\&A}

\subsection{Description and Priority}
Prism AI Browser can extract and summarise web page content, enabling users to quickly understand long articles, technical documents, or e-commerce product pages. This is a \textbf{medium-priority} feature.

\subsection{Implementation Details}
\begin{itemize}[leftmargin=*]
    \item The AI agent executes \texttt{document.body.innerText} via \texttt{Runtime.evaluate} CDP call to extract page text.
    \item The extracted text and user query are assembled into an LLM prompt.
    \item The LLM generates a concise summary or direct answer.
    \item Results are streamed back to the chat UI.
\end{itemize}

\subsection{Functional Requirements Supported}
FR-06

\section{Feature 6: Chromium Rebranding to Prism}

\subsection{Description and Priority}
All user-visible strings in the browser referring to ``Chromium'' are replaced with ``Prism'' to establish a distinct brand identity. This is a \textbf{low-priority} (cosmetic) feature but essential for the final product.

\subsection{Implementation Details}
\begin{itemize}[leftmargin=*]
    \item Modifications are made to \texttt{chrome/app/chromium\_strings.grd}, the central string resource file for Chromium.
    \item Over 100 user-facing strings are updated including: product name, window titles, installer messages, error messages, and taskbar/dock pin prompts.
    \item The short product name becomes ``Prism''; the full description reads: \textit{``Prism is an AI agentic powered web browser that runs webpages and applications with lightning speed.''}.
\end{itemize}

\subsection{Functional Requirements Supported}
FR-01 (brand identity)

\section{Feature 8: Prism Homepage DevSpace}

\subsection{Description and Priority}
The Prism Homepage DevSpace is a custom, developer-oriented new-tab homepage served locally by the Prism browser. It replaces the standard blank new-tab page with a productivity-focused dashboard. This is a \textbf{medium-priority} feature that significantly enhances the day-to-day developer experience.

\subsection{Overview}
On opening a new tab, users are greeted with a full-screen homepage (\texttt{index.html}) featuring a live clock, motivational quotes, animated particle and matrix backgrounds, and a scrollable Dev Space panel. The homepage is designed specifically for the Prism AI Browser but also works in any modern browser when served via a local HTTP server (Python or Node.js).

\subsection{Key Sub-Features}

\subsubsection*{Live Clock and Greetings}
\begin{itemize}[leftmargin=*]
    \item A large, prominent time display (\texttt{\#zen-time}) supports both 12-hour and 24-hour formats selectable from settings.
    \item Personalised motivational greetings are shown as primary and secondary quotes sourced from \texttt{quotes.json} and \texttt{dynamic\_quotes.json}.
    \item The greeting updates daily, creating a fresh experience on each browser session.
\end{itemize}

\subsubsection*{Dev Space Panel}
Scrolling below the clock reveals the Dev Space — a card-based launcher grid offering quick access to developer tools and utilities:
\begin{longtable}{|p{4cm}|p{10cm}|}
\hline
\textbf{Tool Card} & \textbf{Description} \\ \hline
\endhead
Localhost & Quick-launch link to \texttt{http://localhost:3000} (configurable) for local web dev servers \\ \hline
GitHub & Direct link to \texttt{github.com} for version control and repository management \\ \hline
Dev Console & Opens an interactive in-page developer tools panel with browser shortcuts \\ \hline
MDN Docs & Link to Mozilla Developer Network for web API documentation \\ \hline
Color Generator & Opens \texttt{dev-space/color-gen.html} — an interactive colour palette generator \\ \hline
Prompt Synthesizer & Opens \texttt{dev-space/prompt-synthesizer.html} — an AI-powered prompt enhancement and refinement tool \\ \hline
Speed Test & Embedded network speed test via \texttt{dev-space/speedtest/} \\ \hline
Notepad Panel & Floating persistent notepad (\texttt{dev-space/notepad-panel.html}) that saves notes to \texttt{localStorage} \\ \hline
Focus Settings & Configurable focus/distraction-free mode from \texttt{dev-space/focus-settings.html} \\ \hline
Matrix Tester & Visual demo and configurator for the animated Matrix rain background effect \\ \hline
\caption{DevSpace Developer Tools}
\end{longtable}

\subsubsection*{Animated Backgrounds}
\begin{itemize}[leftmargin=*]
    \item \texttt{matrix.js} renders a customisable Matrix-style green rain effect on a canvas element.
    \item Floating particle animations (\texttt{createParticle()}) add ambient depth to the background.
    \item Users can switch between wallpapers from the \texttt{images/Wallpapers/} collection.
\end{itemize}

\subsubsection*{System Information Display}
\begin{itemize}[leftmargin=*]
    \item RAM usage is displayed in real time using the browser's \texttt{performance.memory} API (where available).
    \item Clock style selection is persisted via \texttt{localStorage}, with previews available in the settings panel (\texttt{clock-previews/}).
\end{itemize}

\subsection{Implementation Details}
\begin{itemize}[leftmargin=*]
    \item Built with plain HTML, CSS (\texttt{style.css}, Space Grotesk / JetBrains Mono fonts), and vanilla JavaScript (\texttt{script.js}).
    \item No build step required for the homepage itself; it is served as static files.
    \item Must be served from a local web server (not \texttt{file://}) due to CORS restrictions on \texttt{quotes.json} fetch calls.
    \item Integrates tightly with Prism by linking to internal browser pages via special protocols, with fallback to standard URLs for compatibility.
\end{itemize}

\subsection{Functional Requirements Supported}
Extended browser UX, developer productivity

% ---------------------------------------------------------------

\section{Feature 9: Aura AI Chat (with Image Generation)}

\subsection{Description and Priority}
Aura AI Chat is the primary conversational AI interface integrated into Prism Browser, accessible via a dedicated local web page. It supports multi-model text chat, AI-powered image generation, agentic task execution (browser automation), and persistent conversation history management. This is a \textbf{high-priority} feature representing the core user-facing AI capability of Prism.

\subsection{Overview}
The chat interface (\texttt{index.html}, branded ``Aura -- Your Agentic AI Assistant'') is built as a single-page application with a collapsible sidebar for conversation history, a main chat area, and a bottom input bar supporting both text and voice entry. It uses the \texttt{\#88E755} (Prism green) primary colour scheme and the \textit{Space Grotesk} / \textit{JetBrains Mono} font stack with a dark \texttt{\#0a0a0a} background.

\subsection{AI Models Supported}

\begin{longtable}{|p{3.5cm}|p{3.5cm}|p{7cm}|}
\hline
\textbf{Mode} & \textbf{Model} & \textbf{Description} \\ \hline
\endhead
Aura Mode & Gemini Pro (Google) & Default chat model via Gemini API; best for general conversation, reasoning, and summarisation \\ \hline
DeepSeek V3.2 & DeepSeek-V3.2-Exp (Hugging Face) & Advanced reasoning model; first request may take 20--30 s while the model loads on the Hugging Face inference server \\ \hline
\caption{AI Models Supported by Aura Chat}
\end{longtable}

Users switch between models via the ``Select AI Service'' dropdown in the interface.

\subsection{AI Image Generation}
\begin{itemize}[leftmargin=*]
    \item \textbf{Engine:} FLUX.2 Klein (\texttt{black-forest-labs/flux.2-klein-4b}) via the OpenRouter unified AI gateway API.
    \item \textbf{Trigger:} Users activate image generation by natural language (e.g., ``generate image of a futuristic city'', ``draw a cute robot'').
    \item \textbf{Display:} Generated images are returned as Base64-encoded inline data URIs and rendered directly in the chat. A direct URL fallback is provided.
    \item \textbf{Caching:} Image results are cached to avoid redundant API calls for repeated prompts.
    \item \textbf{Download:} Users can download generated images directly from the chat message.
\end{itemize}

\subsection{Agentic Task Execution}
The \texttt{task-executor.js} module, instantiated as a \texttt{TaskExecutor} class, intercepts commands with recognised patterns using regular expressions and executes real browser actions. Supported task categories:

\begin{longtable}{|p{3.5cm}|p{10cm}|}
\hline
\textbf{Category} & \textbf{Example Commands} \\ \hline
\endhead
YouTube & ``open youtube and search cocomelon'' \\ \hline
Search Engines & ``google search artificial intelligence'', ``search wikipedia quantum physics'' \\ \hline
Social Media & ``open twitter and search AI news'', ``open reddit and search programming'', ``open github and search react'' \\ \hline
Entertainment & ``play imagine dragons on spotify'', ``open netflix and search stranger things'', ``search amazon wireless headphones'' \\ \hline
Productivity & ``set timer for 5 minutes'', ``create event for team meeting tomorrow'', ``check weather in New York'', ``open calculator'' \\ \hline
Navigation & ``open maps coffee shops near me'', ``send email to user@example.com'' \\ \hline
General Web & ``open github.com'', ``open google.com'' \\ \hline
\caption{Supported Agentic Task Execution Categories}
\end{longtable}

All tasks that open external URLs or perform browser navigation require user confirmation via an in-UI dialog. A ``Don't show again'' option is available for frequently approved actions.

\subsection{Conversation History Management}
\begin{itemize}[leftmargin=*]
    \item Conversations are stored persistently in \texttt{localStorage} under named sessions, listed in the collapsible left sidebar.
    \item A search bar (\texttt{\#history-search}) filters conversations by keyword in real time.
    \item Export and import of chat history as JSON files is supported via the sidebar footer buttons.
    \item ``New Chat'' creates a fresh session; ``Clear History'' removes all saved sessions.
\end{itemize}

\subsection{Memory Integration}
Two optional memory backends are supported:
\begin{itemize}[leftmargin=*]
    \item \texttt{local-memory.js}: Stores user context (facts the AI learns about the user) in \texttt{localStorage} for offline, privacy-first persistent memory.
    \item \texttt{supermemory.js}: Optional integration with the Supermemory cloud service for cross-device memory synchronisation (requires API key; opt-in only to preserve privacy goals).
\end{itemize}

\subsection{Error Handling and Reliability}
\begin{itemize}[leftmargin=*]
    \item \textbf{Gemini API 503:} Exponential backoff retry with up to 5 attempts (delays: 1 s, 2 s, 4 s, 8 s, 16 s). Real-time retry status is shown to the user.
    \item \textbf{DeepSeek 503 (model loading):} Same exponential backoff strategy; user is informed that model loading may take 20--30 s.
    \item \textbf{401/403:} Authentication errors with clear user guidance on API key setup.
    \item \textbf{429:} Rate limit detection with a wait-and-retry message.
    \item \textbf{Network errors:} Connection issue detection with actionable error messages.
\end{itemize}

\subsection{Implementation Details}
\begin{itemize}[leftmargin=*]
    \item Built with HTML5, Tailwind CSS (via CLI build -- \texttt{input.css} $\rightarrow$ \texttt{styles.css}), and vanilla JavaScript.
    \item Tailwind custom colour: \texttt{primary: \#88E755} (Prism green); config in \texttt{tailwind.config.js}.
    \item Lucide icons loaded from jsDelivr CDN (switched from unpkg to fix MIME type error).
    \item Production CSS built via \texttt{npm run build:css:prod} (minified); development via \texttt{npm run build:css} (watch mode).
\end{itemize}

\subsection{Functional Requirements Supported}
FR-03 (NLP command processing), FR-04 (voice input), FR-06 (content summarisation), FR-08 (privacy-first local processing)

% ---------------------------------------------------------------

\section{Feature 7: Patch-Based Build and Distribution System}

\subsection{Description and Priority}
Prism maintains all its customisations as a compact set of Git patches applied on top of a pinned Chromium base revision. This allows contributors to reproduce the build on any machine without hosting the full Chromium repository. This is a \textbf{high-priority} operational feature.

\subsection{Implementation Details}
\begin{itemize}[leftmargin=*]
    \item \texttt{BASE\_REVISION.txt}: Records the Chromium commit SHA that Prism patches are based on.
    \item \texttt{patches/0001-prism-working-tree.diff}: Contains all modifications to existing Chromium files (diff format).
    \item \texttt{patches/0002-prism-untracked-new-files.diff}: Contains new files added by Prism (vector icon, AI button source files).
    \item \texttt{scripts/apply\_patches.sh}: Validates the Chromium checkout revision and applies all patches in order. Supports options \texttt{--check}, \texttt{--skip-base-check}, and \texttt{--patch-dir}.
    \item \texttt{scripts/build\_prism.sh}: One-command wrapper that applies patches and invokes \texttt{gn gen} + \texttt{ninja} to build the \texttt{chrome} target.
\end{itemize}

\needspace{8cm}
\subsection{Functional Requirements Supported}
The patch-based build and distribution system is a foundational infrastructure component that underpins every functional requirement in this document.
Without a reproducible, patch-managed build pipeline, none of the Prism-specific features could be reliably compiled, tested, or delivered to end users.
Specifically:

\begin{itemize}[leftmargin=*]
    \item \textbf{FR-01 -- FR-08 (all features):} Every Prism capability -- the AI toolbar button, CDP integration, AI agent server, voice processing, tab management, form automation, content summarisation, and privacy-first inference -- is delivered exclusively through the patch set applied on top of Chromium. A broken or non-reproducible patch application would prevent any of these features from functioning.
    \item \textbf{Reproducible Builds:} \texttt{BASE\_REVISION.txt} pins the exact Chromium commit SHA, ensuring that the C++ patches in \texttt{patches/0001-*} and \texttt{patches/0002-*} always apply cleanly against the intended source tree, regardless of upstream Chromium development activity.
    \item \textbf{Complete Runtime Deployment:} The packaging stage (Stage~6) bundles the compiled Prism binary together with the AI agent (\texttt{ai-chat/}) and the custom new-tab homepage (\texttt{homepage-devspace/}) into a single distributable archive. This ensures all runtime components required by FR-01 through FR-08 are co-deployed and version-matched.
    \item \textbf{Patch Drift Detection:} If an upstream Chromium change modifies a file touched by a Prism patch, \texttt{apply\_patches.sh} fails immediately with a conflict report, protecting the integrity of all features that depend on those patched files.
\end{itemize}
